\subsection{Context}

We are now well aware that we need to significantly reduce our greenhouse gas emissions to limit global warming. Electrifying as many sectors as possible and decarbonizing electricity production are two major ways to achieve this goal. However, the renewable energy sources used to produce clean electricity, such as solar and wind power, are intermittent and require much more space than traditional power plants. This has led to growing interest in decentralized electricity production, as it can spread generation over a larger area to reduce intermittency issues and take advantage of the many available surfaces. For example, urban areas have numerous unused surfaces that could be used to install solar panels, such as rooftops and building facades.

By installing renewable energy sources at home or in local areas—particularly solar panels—a former consumer can become a “prosumer,” meaning they both consume and produce electricity within the grid. We can then distinguish two cases of production: either the prosumer injects power into the grid, or the prosumer self-consumes the produced power. Depending on policies and grid fees, self-consumption is often much more attractive for prosumers than injecting power into the grid.

To maximize self-consumption, we can generally consider two main approaches. The first is to install energy storage systems to store excess production for later use. However, storage systems are still quite expensive. The second approach is to adapt consumption to production, which can be difficult to achieve at the individual level. This is where local energy communities (LECs) become useful. By grouping prosumers together, they can share production and consumption, thereby maximizing self-consumption at the community level. Of course, LECs can also benefit from storage systems. When a community is built using only renewable energy sources and storage systems, we refer to it as a “renewable energy community” (REC).

This study will focus on LECs and, more specifically, on RECs in Europe, as this concept depends heavily on local regulations. The results will still probably remain valid in different regions of the world, but no extensive study on this subject were done as part of this work.

Energy communities were introduced in the EU in 2016 under the Clean Energy Package, particularly through the Renewable Energy Directive, and have since grown significantly. The adoption of legislative measures and the Directive on Common Rules for the Internal Electricity Market provided the first formal definitions of energy communities in European law. RECs typically connect to public energy networks operated by local DSOs, but the 2019/944/CE Directive promotes closed distribution systems (CDSs). These are intended for the exclusive use of REC members and are expected to organize exchanges within the REC. The European directive does not impose restrictions on how RECs may distribute profits among participants, so this responsibility would also fall to the CDSs.

In most European countries, energy communities are limited in geographic scope and number of participants, based on the principle that they should be connected to the same feeder. Nevertheless, there is considerable diversity in the size of energy communities, ranging from a few households at low voltage to industrial parks at medium voltage.

\subsection{Problem statement}

For maximizing self-consumption, RECs need to be operated in a way that optimizes the use of the available renewable energy. This can be achieved by implementing a local energy management system (EMS) that coordinates production, consumption and storage within the community. However, this operation can be done in many ways, each of them raising different type of concern. Among them, we can mention, the data privacy of the members of the community, the computational complexity of the EMS, the robustness of the system to failures, or the fairness of the gain allocation among the members of the community. 

In this study we will focus on this last point. When a community is built, the members of the community need to agree on how to share the gains resulting from the operation of the community. This is a crucial point for a community to be successful, as it can lead to conflicts and dissatisfaction among the members if not handled properly. Therefore, it is important to find a fair way to allocate the gains among the members of the community. Of course fairness can be defined in many ways, and should take into account different criteria. In \cite{these_sorel}, the author uses a sociological profile of the members to define what the members want from the community. We will then need to evaluate the different gain allocation functions that can be used.

But in the different studies on this topic, it is often hard to understand and visualize what are the impacts of the different gain allocation functions on the fairness within the community. This will be the main objective of this study. First we will model an energy community and try to visualize the effect of the different way of operating the community. We will try to understand the effect of the sociological profile as well as the impact of the different gain allocation functions on the fairness within the community.

\subsection{Overview of the report}

In the first part of this report, we present the state of the art on energy community modeling and the associated fairness challenges. We then describe the methodology used to conduct the study. Next, we introduce the mathematical model developed to represent the energy community. The gain allocation functions evaluated in this work are subsequently detailed. Finally, we employ various visualization tools to interpret and discuss the simulation results.